% 来源：OEBPS/Text/01_29.xhtml（OPF spine 顺序；仅第一册）
\cleardoublepage
\BookEntry[大卫·李嘉图《政治经济学及赋税原理》]{大卫·李嘉图《政治经济学及赋税原理》}

\BookCover{images/书单15政治经济学及赋税原理.jpg}{[英] 大卫·李嘉图著| 郭大力、王亚南译| 商务印书馆2013年版}

\begin{BookSummaryBox}
\textbf{英国古典政治经济学的巅峰之作}
\par
\end{BookSummaryBox}

\subsubsection*{这本书谈什么？}

《政治经济学及赋税原理》批判继承了亚当·斯密的劳动价值论，并以此为基础阐述了工资、利润和地租的相互关系，以及货币理论、对外贸易中的比较成本学说、赋税的一般原理和原则。
\par

在李嘉图看来，人类行为的准则是个人利益，是追求个人功利，即寻求快乐和避免痛苦，而追求个人利益同整个社会的利益是一致的，符合社会最大多数人的最大幸福。他实际上是站在工业资产阶级立场上，把资本家追求利润的个人利益与生产发展、社会财富增加和社会进步看作是一致的。不过，这种研究方法带有形而上学的片面性，把资本主义生产关系和制度看作是永恒的、不变的。
\par

\subsubsection*{芒格为什么谈起这本书？}

2003年10月3日，芒格在“论学院派经济学”的演讲中谈道：
\par

\begin{quotation}
如果美国和中国进行贸易，那么美国人的生活水平会得到改善，李嘉图证明了这一点。但是哪个国家的经济会发展得更快呢？那显然是中国。在自由贸易的推动之下，他们吸收了世界上各种现代科技，而且正如亚洲四小龙已经证明的，他们很快就会走到前面去。看看中国的香港和台湾，看看早年的日本就知道了。
\par
\end{quotation}

比较优势通过自由贸易带来的好处通常遭到忽略，在其1817年出版的名著《政治经济学及赋税原理》中，大卫·李嘉图对此进行了揭示：
\par

\begin{quotation}
葡萄牙生产红酒和毛呢的成本可能比英国低。然而，这两个国家生产这两种产品的相对成本是有差别的。在英国，生产红酒很难，但生产毛呢的难度并不是很大。在葡萄牙，这两者都很容易生产。因此，尽管葡萄牙生产毛呢的成本比英国低，但如果葡萄牙生产更多的红酒，用来交换英国的毛呢，那么葡萄牙的毛呢成本会更低。相反地，英国也从这种贸易中获益，因为它生产毛呢的成本没有变化，但它现在能够用更接近于毛呢的成本来获得红酒。
\par
\end{quotation}

芒格同时还指出：
\par

\begin{quotation}
国家之间相互“委派”任务能够获得比较优势，同样地，管理人员在安排工作时也可以利用李嘉图的这个原理，而这是很多人都没想到的。即使一个管理人员能够亲自更好地完成各种工作，把这些工作分散到各人手里仍然是对大家都有利的做法。
\par
\end{quotation}

\subsubsection*{《政治经济学及赋税原理》作者大卫·李嘉图如是说}

“金属和贵金属同其他物品一样，是由劳动生产的，因而其价值也是由劳动决定的，并且是随着劳动量的变动而变动的。”
\par

“货币的功能主要是流通的手段。”
\par

“当流通中的货币数量多于正常需要水平时，商品价格就上涨；相反，当流通中货币数量小于正常需要水平时，商品价格下跌。假定一国的商品流通量和价值是已知数，则‘一国所能投入的货币量必须取决于货币的价值’。”
\par
